Prime gaps are often studied through their sizes, asymptotic behavior, and distributional regularities. A complementary finite-computational view is to track the residue classes of consecutive primes after removing the smallest prime divisibility constraints. For primes greater than 5, residues modulo 30 lie in
\[
R=\{1,7,11,13,17,19,23,29\}.
\]
This gives an eight-state transition system. A first-order Markov model is a natural baseline because it records the observed probabilities of moving from one residue class to the next. However, first-order transitions need not capture memory that appears over two or more steps. This paper asks whether the empirical two-step residue transition operator contains structure beyond the first-order Markov prediction.

The central diagnostic is the residual operator
\[
\Delta=P^{(2)}-P^2,
\]
where \(P^{(2)}\) is the empirical two-step transition operator and \(P^2\) is the two-step prediction implied by the first-order operator. If the residue process were adequately described by \(P\), this residual would be small and structureless up to finite-sample effects. Instead, the notebooks summarized here identify stable low-rank components in \(\Delta\), validate them against synthetic controls, and package the results for reproducible inspection.
